%-*- coding:UTF-8 -*-
%作者：白鸽坐飞机（知乎、简书和B站同号）
% SYSUplate.tex
% 中山大学博士毕业论文模板
% 请使用 XeLaTeX编译本文。也可以选择XeLaTeX+MakeIndex+BibTex编译



\documentclass[UTF8,12pt,twoside,a4paper]{ctexart} %twoside 双面排版，奇偶页的页眉页脚和页边距可以不同

%应用的包——————————————————————————————————————————
\usepackage{fontspec}%字体的包
\newcommand{\tnewroman}{\fontspec{Times New Roman}}
\newcommand{\adobegothicstdb}{\fontspec{Adobe Gothic Std B}}
\newcommand{\Centaur}{\fontspec{Centaur}}
\setmainfont{Times New Roman}%fontspec下这个命令设置全局默认字体


%页眉参考http://www.ctex.org/documents/packages/layout/fancyhdr.htm

\usepackage{graphicx}%插图宏包
\usepackage{caption}%figure legend

\usepackage{setspace}%设置行距的包
%\setlength{\baselineskip}{1.6em}%2倍行距


\usepackage{indentfirst}%首行缩进的包
%\setlength{\parindent}{2em}%首行缩进2

\usepackage{geometry}%设置页边距的包
\geometry{left=2.5cm,right=2.5cm,top=3.5cm,bottom=2.5cm} % 设置页边距


\usepackage{titletoc} %更改目录的包
\titlecontents{section}[3em]%标题位置左间距
{\songti \zihao{-4}\vspace{10pt}}           %设置标题字体是宋体（但是被后面命令覆盖了）；字号小四；与上一个标题垂直距离
{\contentslabel{2em}}%设置标题标志的格式，如序号格式、序号宽度、序号与标题内容之间的间距
{\hspace*{-2em}}%设置无序号标题的格式，如字体、字体尺寸、位置等。该参数可以空置
{~\titlerule*[0.6pc]{$.$}~\contentspage}%设置标题与页码之间的指引线样式以及页码的格式，该参数如果空置，标题将无指引线

\titlecontents{subsection}[4em]{\songti\zihao{-4}\vspace{10pt}}{\contentslabel{2em}}{\hspace*{-2em}}{~\titlerule*[0.6pc]{$.$}~\contentspage} %指引线与页码

\usepackage{enumerate} %设置列表环境的包

\bibliographystyle{unsrt} %声明参考文献格式为plain,也可以选unsrt,alpha,abbrv等

\usepackage{fancyhdr}  %设置页眉页码的包
\pagestyle{fancy}
\lhead{中山大学博士学位论文} %页眉左边
\chead{} %页眉中间
\rhead{博士毕业论文题目比如三国演义某某某} %页眉右边
%\lfoot{From: K. Grant} %页脚右边是From: K. Grant
\cfoot{\thepage} %页脚中间是本页页码
%\rfoot{To XXX} %页脚右边是To XXX
\renewcommand{\headrulewidth}{0.4pt} %页眉划线
%\renewcommand{\footrulewidth}{0.4pt}%页脚划线

%\pagestyle{myheadings}
%\markboth{中山大学博士学位论文}{战猇亭先主得仇人~守江口书生拜大将}%设置左边页眉的中山大学博士学位论文，右边页眉是论文标题



\usepackage{afterpage}%更改目录计数的

\usepackage{hyperref}%内部交叉引用








\begin{document}


%封面________________________________________________________________
	\begin{titlepage}
	
		\begin{center}
			\includegraphics{SYSU.pdf}\\%插入pdf的第1页	 {graphicx}包内的命令
			\heiti{\zihao{0}博士学位论文}\\ %heiti就是黑体字。0就是初号字。\\表示另起一行。
			\vspace*{36pt} %行间距增加36pt
			\heiti{\zihao{2}
博士毕业论文题目比如三国演义某某某
}\\%\zihao{2}就是二号字；-5就是小五 ;~表示此处有空格，但无需换行
			\vspace*{16pt}%行间距增加16pt
			\heiti{\zihao{2}
You English Title for the PhD Degree Such as Story of Three Country
}\\
			\vspace*{36pt}%行间距增加36pt
			\songti
			\zihao{-3}
			\textbf{
				{专 业：\quad 战斗学}\\ %\quad表示空格
				{博士研究生：\quad 陆逊}\\
				{指导老师：\quad 孙权}\\
						}
				\vspace*{36pt} %行间距增加36pt
		\end{center}	
 		
		\begin{flushleft}     % 左对齐    
\zihao{-3}
答辩委员会主席：\\
答辩委员会成员：\\
		\end{flushleft}

		\begin{center}
		\vspace*{90pt} %行间距增加90pt
		\zihao{-3}
      中山大学下面某个具体的学院\\
		2020年6月\quad \quad 广州
		\end{center}

	\end{titlepage}
\clearpage{\pagestyle{empty}\cleardoublepage}%空余一页，余下内容直接进入下一页，而且不要出现页眉页码
%\cleardoublepage %这一页留空，余下内容直接进入新的一页
%\clearpage 

%原创性声明与授权书————————————————————————————————————
	%\begin{titlepage}
		
		\begin{center}
		\heiti{\zihao{-3}论文原创性声明}\\
		\end{center}

	\begin{flushleft}%左对齐开始     flushright则是右对齐
						\songti
                  \zihao{4}
						 \renewcommand{\baselineskip}{1.6em}	%双倍行距
						
\quad\quad 本人郑重声明：所呈交的学位论文，是本人在导师的指导下，独立进行研究工作所取得的成果。除文中已经注明引用的内容外，本论文不包含任何其他个人或集体已经发表或撰写过的作品成果。对本文的研究作出重要贡献的个人和集体，均已在文中以明确方式标明。本人完全意识到本声明的法律结果由本人承担。
	\end{flushleft}

	\begin{flushright}     % 左对齐    
	\zihao{4}
	\renewcommand{\baselineskip}{1.6em}	%双倍行距
学位论文作者签名：\quad\quad\quad\quad\quad\quad \\
日期：\quad\quad 年  \quad\quad  月  \quad\quad  日\quad\quad\quad\quad \\
		\end{flushright}

	\begin{center}
	\vspace*{56pt} %行间距增加42pt
	\heiti{\zihao{-3}学位论文使用授权声明}\\
		\end{center}


	\begin{flushleft}%左对齐开始     flushright则是右对齐
						\songti
                  \zihao{4}
						 \renewcommand{\baselineskip}{1.6em}	%双倍行距
						
\quad\quad 本人完全了解中山大学有关保留、使用学位论文的规定，即：学校有权保留学位论文并向国家主管部门或其指定机构送交论文的电子版和纸质版；有权将学位论文用于非盈利目的的少量复制并允许论文进入学校图书馆、院系资料室被查阅；有权将学位论文的内容编入有关数据库进行检索；可以采用复印、缩印或其他方法保存学位论文；可以为存在馆际合作关系的兄弟高校用户提供文献传递服务和交换服务。\\
\quad\quad 保密论文保密期满后，适用本声明。\\

	\end{flushleft}
	
	\begin{flushleft}     % 左对齐    
\zihao{4}
\renewcommand{\baselineskip}{1.6em}	%双倍行距
学位论文作者签名：\quad\quad\quad\quad\quad\quad 导师签名：\\
日期：\quad\quad 年  \quad\quad  月  \quad\quad  日\quad\quad\quad  日期：\quad\quad 年  \quad\quad  月  \quad\quad  日\quad\quad\quad\quad\\
		\end{flushleft}
%	\end{titlepage}
%\cleardoublepage %这一页留空，余下内容直接进入新的一页
\clearpage{\pagestyle{empty}\cleardoublepage}%空余一页，余下内容直接进入下一页，而且不要出现页眉页码

%______________________________________________________________________________________________





\setcounter{page}{1}%页码开始计数

\pagenumbering{Roman}%样式有:  arabic, roman, Roman, alph, Alph





%中文摘要——————————————————————————————————————
\begin{titlepage}

		
		\addcontentsline{toc}{section}{\songti{\zihao{-4}\textbf{中文摘要}}}
		\begin{center}
			\zihao{-1}\bfseries 
博士毕业论文题目比如三国演义某某某
\\%\zihao{-2}字号-4就是小四；-5就是小五 
			\vspace*{36pt}%行间距增加36pt
			\heiti
			\zihao{-3}
				{专\quad \quad 业：\quad 战斗学}\\
				{博士研究生：\quad 陆逊}\\
				{指导老师：\quad 孙权}\\
		 \vspace*{36pt}
%\section*{\heiti{\zihao{-2}{中文摘要}}}
		\heiti{\zihao{-2}{中文摘要}}
				
				\vspace*{16pt}
		\end{center}
		
						%\indent\setlength{\parindent}{2em}%首行缩进2字符
						\songti
                  \zihao{-4}
						 %\renewcommand{\baselineskip}{1.6em}
							%\hspace{0.8em}%首段首行缩进
却说章武二年春正月，武威后将军黄忠随先主伐吴；忽闻先主言老将无用，即提刀上马，引亲随五六人，径到彝陵营中。吴班与张南、冯习接入，问曰：“老将军此来，有何事故？”忠曰：“吾自长沙跟天子到今，多负勤劳。今虽七旬有余，尚食肉十斤，臂开二石之弓，能乘千里之马，未足为老。昨日主上言吾等老迈无用，故来此与东吴交锋，看吾斩将，老也不老！”正言间，忽报吴兵前部已到，哨马临营。忠奋然而起，出帐上马。冯习等劝曰：“老将军且休轻进。”忠不听，纵马而去。吴班令冯习引兵助战。忠在吴军阵前，勒马横刀，单搦先锋潘璋交战。璋引部将史迹出马。迹欺忠年老，挺枪出战；斗不三合，被忠一刀斩于马下。潘璋大怒，挥关公使的青龙刀，来战黄忠。交马数合，不分胜负。忠奋力恶战，璋料敌不过，拨马便走。

忠乘势追杀，全胜而回。路逢关兴、张苞。兴曰：“我等奉圣旨来助老将军；既已立了功，速请回营。”忠不听。次日，潘璋又来搦战。黄忠奋然上马。兴、苞二人要助战，忠不从；吴班要助战，忠亦不从；只自引五千军出迎。战不数合，璋拖刀便走。忠纵马追之，厉声大叫曰：“贼将休走！吾今为关公报仇！”追至三十余里，四面喊声大震，伏兵齐出：右边周泰，左边韩当，前有潘璋，后有凌统，把黄忠困在垓心。忽然狂风大起，忠急退时，山坡上马忠引一军出，一箭射中黄忠肩窝，险些儿落马。吴兵见忠中箭，一齐来攻，忽后面喊声大起，两路军杀来，吴兵溃散，救出黄忠，乃关兴、张苞也。二小将保送黄忠径到御前营中。忠年老血衰，箭疮痛裂，病甚沉重。先主御驾自来看视，抚其背曰：“令老将军中伤，朕之过也！”忠曰：“臣乃一武夫耳，幸遇陛下。臣今年七十有五，寿亦足矣。望陛下善保龙体，以图中原！”言讫，不省人事。是夜殒于御营。后人有诗叹曰：“老将说黄忠，收川立大功。重披金锁甲，双挽铁胎弓。胆气惊河北，威名镇蜀中。临亡头似雪，犹自显英雄。”
		
		
%\heiti{关键词：刘备；陆逊；书生；大将}\\

\begin{flushleft}
\heiti{\zihao{3}{目的：}}
\end{flushleft}

\begin{enumerate}%有序列表。无序列表itemize。描述列表description
	\songti
			\zihao{-4}
  \item 滚滚长江东逝水，浪花淘尽英雄。
  \item 是非成败转头空。
  \item 青山依旧在，几度夕阳红。
\end{enumerate}


\begin{flushleft}
\heiti{\zihao{3}{方法：}}
\end{flushleft}

\begin{enumerate}%有序列表。无序列表itemize。描述列表description
	\songti
			\zihao{-4}
  \item 白发渔樵江渚上，惯看秋月春风。
  \item 一壶浊酒喜相逢。
  \item 忠乘势追杀，全胜而回。
	\item 古今多少事，都付笑谈中。
  \item 话说天下大势，分久必合，合久必分。
  \item 周末七国分争，并入于秦。
\end{enumerate}


\begin{flushleft}
\heiti{\zihao{3}{结果：}}
\end{flushleft}



1.\heiti  及秦灭之后，楚、汉分争，又并入于汉。\\
\songti 汉朝自高祖斩白蛇而起义，一统天下，后来光武中兴，传至献帝，遂分
为三国。

2.\heiti 推其致乱之由，殆始于桓、灵二帝。\\
\songti 桓帝禁锢善类，崇信宦官。及桓帝崩，灵帝即
位，大将军窦武、太傅陈蕃共相辅佐。

3.\heiti  时有宦官曹节等弄权，窦武、陈蕃谋诛之，机事不
密，反为所害，中涓自此愈横。\\
\songti 建宁二年四月望日，帝御温德殿。方升座，殿角狂风骤起。只见一条大青蛇，从梁上飞
将下来，蟠于椅上。

4.\heiti 帝惊倒，左右急救入宫，百官俱奔避。
\songti 须臾，蛇不见了。忽然大雷大
雨，加以冰雹，落到半夜方止，坏却房屋无数。建宁四年二月，洛阳地震；又海水泛溢，沿
海居民，尽被大浪卷入海中。
1.\heiti  光和元年，雌鸡化雄。\\
\songti   六月朔，黑气十余丈，飞入温德殿中。
秋七月，有虹现于玉堂；五原山岸，尽皆崩裂。种种不祥，非止一端。帝下诏问群臣以灾异
之由，议郎蔡邕上疏，以为蜺堕鸡化，乃妇寺干政之所致，言颇切直。

5.\heiti 帝览奏叹息，因起更衣。\\
\songti 曹节在后窃视，悉宣告左右；遂以他事陷邕于罪，放归田里。后张让、赵忠、封谞、段
珪、曹节、侯览、蹇硕、程旷、夏恽、郭胜十人朋比为奸，号为“十常侍”。


\begin{flushleft}
\heiti{\zihao{3}{结论：}}
\end{flushleft}

\begin{enumerate}%有序列表。无序列表itemize。描述列表description
	\songti
			\zihao{-4}
  \item 帝尊信张让，呼为“阿父”。朝政日非，以致天下人心思乱，盗贼蜂起。
  \item 时巨鹿郡有兄弟三人，一名张角，一名张宝，一名张梁。
  \item  那张角本是个不第秀才，因入山采药，遇一老人，碧眼童颜，手执藜杖，
	\item 唤角至一洞中，以天书三卷授之，曰：“此名《太平要术》，汝得之，当代天宣化，普救世人；若萌异心，必获恶报。”
  \item  角拜问姓名。
  \item  老人曰：“吾乃南华老仙也。”
\end{enumerate}

\begin{flushleft}
\heiti{\zihao{3}{关键词：}} \songti{\zihao{4}{黄忠，关兴，张苞}}
\end{flushleft}

		%\end{titlepage}
\clearpage 
%\cleardoublepage %这一页留空，余下内容直接进入新的一页




%英文摘要——————————————————————————————————————
%\begin{titlepage}

\begin{center}
\addcontentsline{toc}{section}{{\setmainfont{Times New Roman}}\zihao{-4}{Abstract}}  
			\zihao{-1} \bfseries
    Liu Bei killed the enemy in Xiaoting and Lu Xvn became a general in Jiangkou
        \\
			\vspace*{36pt}%行间距增加36pt
			\setmainfont{Times New Roman}
			\zihao{3}
			\textbf{
				{Major：\quad Battle}\\
				{Candidate：\quad Lu Xvn}\\
				{Supervisor：\quad Sun Quan}}\\
				\vspace*{36pt}

 {\setmainfont{Times New Roman}}
	\bfseries{\zihao{2}{Abstract}}
\end{center}
%\hspace{0.2em}%首段首行缩进

Teams from Britain and the EU are continuing efforts to break the deadlock over agreements on future relations. The transition period ends in just over two weeks, and two sides haven't agreed rules for fair competition, the settling of disputes and fishing rights.

%\indent\setlength{\parindent}{1em}%首行缩进2字符
The Head of the Nigerian military has inspected the operation to free children abducted from a school in north on Friday. Officials said the gunmen are believed to hold ten children hostage. But staff of the school believe this is underestimated.

Officials in Niger say at least 27 people have been killed in what they described as a savage attack on a southeastern village. They believe the Islamic Militants are seeking to intimidate voters before local elections.

The Head of the Nigerian military has inspected the operation to free children abducted from a school in north on Friday. Officials said the gunmen are believed to hold ten children hostage. But staff of the school believe this is underestimated.

Officials in Niger say at least 27 people have been killed in what they described as a savage attack on a southeastern village. They believe the Islamic Militants are seeking to intimidate voters before local elections.


\begin{flushleft}
{\setmainfont{Times New Roman}}
	\bfseries{\zihao{3}{Objective:}}
%\heiti{\zihao{3}{Objective:}}
\end{flushleft}

\begin{enumerate}%有序列表。无序列表itemize。描述列表description
			\zihao{-4}
  \item A new ceasefire has come into effect in the disputed region of Nagorno-Karabakh. 
  \item  The humanitarian truce was announced after the Russian foreign minister Sergey Lavrov spoke to both Armenia and Azerbaijan. President Macron of France said the new ceasefire must be strictly respected by both sides.
  \item Two weeks before the presidential election in Ivory Coast, clashes have broken out in the east.
\end{enumerate}


\begin{flushleft}
{\setmainfont{Times New Roman}}
	\bfseries{\zihao{3}{Methods:}}
%\heiti{\zihao{3}{Methods:}}
\end{flushleft}

\begin{enumerate}%有序列表。无序列表itemize。描述列表description
				\zihao{-4}
  \item A new ceasefire has come into effect in the disputed region of Nagorno-Karabakh. 
  \item  The humanitarian truce was announced after the Russian foreign minister Sergey Lavrov spoke to both Armenia and Azerbaijan. President Macron of France said the new ceasefire must be strictly respected by both sides.
  \item Two weeks before the presidential election in Ivory Coast, clashes have broken out in the east.
	\item A new ceasefire has come into effect in the disputed region of Nagorno-Karabakh. 
  \item  The humanitarian truce was announced after the Russian foreign minister Sergey Lavrov spoke to both Armenia and Azerbaijan. President Macron of France said the new ceasefire must be strictly respected by both sides.
  \item Two weeks before the presidential election in Ivory Coast, clashes have broken out in the east.
	\item A new ceasefire has come into effect in the disputed region of Nagorno-Karabakh. 
  \item  The humanitarian truce was announced after the Russian foreign minister Sergey Lavrov spoke to both Armenia and Azerbaijan. President Macron of France said the new ceasefire must be strictly respected by both sides.
\end{enumerate}

\begin{flushleft}
{\setmainfont{Times New Roman}}
	\bfseries{\zihao{3}{Results:}}
\end{flushleft}

\bfseries{1.A new ceasefire has come into effect in the disputed region of Nagorno-Karabakh.}\\
 \mdseries The humanitarian truce was announced after the Russian foreign minister Sergey Lavrov spoke to both Armenia and Azerbaijan.

\bfseries{2.A new ceasefire has come into effect in the disputed region of Nagorno-Karabakh.}\\
 \mdseries The humanitarian truce was announced after the Russian foreign minister Sergey Lavrov spoke to both Armenia and Azerbaijan.

\bfseries{3.A new ceasefire has come into effect in the disputed region of Nagorno-Karabakh.}\\
 \mdseries The humanitarian truce was announced after the Russian foreign minister Sergey Lavrov spoke to both Armenia and Azerbaijan.

\bfseries{4.A new ceasefire has come into effect in the disputed region of Nagorno-Karabakh.}\\
 \mdseries The humanitarian truce was announced after the Russian foreign minister Sergey Lavrov spoke to both Armenia and Azerbaijan.

\begin{flushleft}
{\setmainfont{Times New Roman}}
	\bfseries{\zihao{3}{Conclusion:}}
%\heiti{\zihao{3}{Objective:}}
\end{flushleft}

\begin{enumerate}%有序列表。无序列表itemize。描述列表description
			\zihao{-4}
  \item A new ceasefire has come into effect in the disputed region of Nagorno-Karabakh. 
  \item  The humanitarian truce was announced after the Russian foreign minister Sergey Lavrov spoke to both Armenia and Azerbaijan. President Macron of France said the new ceasefire must be strictly respected by both sides.
  \item Two weeks before the presidential election in Ivory Coast, clashes have broken out in the east.
\end{enumerate}


\begin{flushleft}
{\setmainfont{Times New Roman}}
	\bfseries{\zihao{3}{Key Words:}}
%\heiti{\zihao{3}{Objective:}}
\end{flushleft}

LiuBei; luXvn; SunQuan;YiLingzhizhan

\end{titlepage}

%\cleardoublepage %这一页留空，余下内容直接进入新的一页
\clearpage{\pagestyle{empty}\cleardoublepage}%空白页不要页眉页码


%目录————————————————————————————————————————————
\tableofcontents%形成目录
\clearpage{\pagestyle{empty}\cleardoublepage}
%\clearpage %这一页留空，余下内容直接进入新的一页
%\newpage%插入一个完整的没有内容的空白页
\mbox{}
%\cleardoublepage %这一页留空，余下内容直接进入新的一页


%—第一章~前言———————————————————————————————————————
\setcounter{page}{1}
\pagenumbering{arabic}%阿拉伯数字页码

\addcontentsline{toc}{section}{\textbf{\songti{\zihao{-4}{\bfseries 第一章~前言}}}}
	\begin{center}
	\heiti{\zihao{2}{第一章~前言}}
	\end{center}
\vspace*{10pt}

	


	
	\begin{flushleft}%左对齐开始     flushright则是右对齐

\heiti{\zihao{3}{
小标题1 	宴桃园豪杰三结义， 斩黄巾英雄首立功
}}					

\indent\setlength{\parindent}{2em}%首行缩进2字符
						\songti
                  %\zihao{-4}
						 %\renewcommand{\baselineskip}{1.5em}
						\hspace{1.6em}%首段首行缩进
却说韩当、周泰探知先主移营就凉，急来报知陆逊。逊大喜，遂引兵自来观看动静；只
见平地一屯，不满万余人，大半皆是老弱之众，大书“先锋吴班”旗号。周泰曰：“吾视此
等兵如儿戏耳。愿同韩将军分两路击之。如其不胜，甘当军令。”陆逊看了良久，以鞭指
曰：“前面山谷中。隐隐有杀气起；其下必有伏兵，故于平地设此弱兵，以诱我耳。诸公切
不可出。”众将听了，皆以为懦。

次日，吴班引兵到关前搦战，耀武扬威，辱骂不绝；多有解衣卸甲，赤身裸体，或睡或
坐。徐盛、丁奉入帐禀陆逊曰：“蜀兵欺我太甚！某等愿出击之！”逊笑曰：“公等但恃血
气之勇，未知孙、吴妙法，此彼诱敌之计也：三日后必见其诈矣。”徐盛曰：“三日后，彼
移营已定，安能击之乎？”逊曰：“吾正欲令彼移营也。”诸将哂笑而退。过三日后，会诸
将于关上观望，见吴班兵已退去。逊指曰：“杀气起矣。刘备必从山谷中出也。”言未毕，
只见蜀兵皆全装惯束，拥先主而过。吴兵见了，尽皆胆裂。逊曰：“吾之不听诸公击班者，
正为此也。今伏兵已出，旬日之内，必破蜀矣。”诸将皆曰：“破蜀当在初时，今连营五六
百里，相守经七八月，其诸要害，皆已固守，安能破乎？”逊曰：“诸公不知兵法。备乃世
之枭雄，更多智谋，其兵始集，法度精专；今守之久矣，不得我便，兵疲意阻，取之正在今
日。”诸将方才叹服。后人有诗赞曰：“虎帐谈兵按六韬，安排香饵钓鲸鳌。三分自是多英
俊，又显江南陆逊高。”却说陆逊已定了破蜀之策，遂修笺遣使奏闻孙权，言指日可以破蜀
之意。权览毕，大喜曰：“江东复有此异人，孤何忧哉！诸将皆上书言其懦，孤独不信，今
观其言，果非懦也。”于是大起吴兵来接应。却说先主于猇亭尽驱水军，顺流而下，沿江屯
扎水寨，深入吴境。黄权谏曰：“水军沿江而下，进则易，退则难。臣愿为前驱。陛下宜在
后阵，庶万无一失。”先主曰：“吴贼胆落，朕长驱大进，有何碍乎？”众官苦谏，先主不
从。遂分兵两路：命黄权督江北之兵，以防魏寇；先主自督江南诸军，夹江分立营寨，以图
进取。细作探知，连夜报知魏主，言蜀兵伐吴，树栅连营，纵横七百余里，分四十余屯，皆
傍山林下寨；今黄权督兵在江北岸，每日出哨百余里，不知何意。魏主闻之，仰面笑曰：
“刘备将败矣！”群臣请问其故。魏主曰：“刘玄德不晓兵法；岂有连营七百里，而可以拒
敌者乎？包原隰险阻屯兵者，此兵法之大忌也。\cite{RN1}

	\end{flushleft}%左对齐结束


	\begin{flushleft}
\heiti{\zihao{3}{
小标题2	张翼德怒鞭督邮， 何国舅谋诛宦竖
}}				

\indent\setlength{\parindent}{2em}%首行缩进2字符
						\songti
                  %\zihao{-4}
						% \renewcommand{\baselineskip}{1.5em}
						\hspace{1.6em}%首段首行缩进
玄德必败于东吴陆逊之手，旬日之内，消息必至矣。”群臣犹未信，皆请拨兵备之。魏主曰：
“陆逊若胜，必尽举吴兵去取西川；吴兵远去，国中空虚，朕虚托以兵助战，令三路一齐进兵，
东吴唾手可取也。”众皆拜服。魏主下令，使曹仁督一军出濡须，曹休督一军出洞口，
曹真督一军出南郡：“三路军马会合日期，暗袭东吴。朕随后自来接应。”调遣已定。不说魏兵袭吴。且说马良至川，入见孔明，呈上图本而言曰：“今移营夹江，横占七百里，下四十余屯，皆依溪傍涧，林木茂盛之处。皇上令良将图本来与丞相观之。”孔明看讫，拍案叫苦曰：“是何人教主上如此下寨？可斩
此人！”\cite{RN4}马良曰：“皆主上自为，非他人之谋。”孔明叹曰：“汉朝气数休矣！”良问其
故。孔明曰：“包原隰险阻而结营，此兵家之大忌。倘彼用火攻，何以解救？又，岂有连营
七百里而可拒敌乎？祸不远矣！陆逊拒守不出，正为此也。汝当速去见天子，改屯诸营，不
可如此。”良曰：“倘今吴兵已胜，如之奈何？”孔明曰：“陆逊不敢来追，成都可保无
虞。”良曰：“逊何故不追？”孔明曰：“恐魏兵袭其后也。主上若有失，当投白帝城避
之。吾入川时，已伏下十万兵在鱼腹浦矣。”良大惊曰：“某于鱼腹浦往来数次，未尝见一
卒，丞相何作此诈语？”孔明曰：“后来必见，不劳多问。”马良求了表章，火速投御营
来。孔明自回成都，调拨军马救应。却说陆逊见蜀兵懈怠，不复提防，升帐聚大小将士听令
曰：“吾自受命以来，未尝出战。今观蜀兵，足知动静，故欲先取江南岸一营。谁敢去
取？”言未毕，韩当、周泰、凌统等应声而出曰：“某等愿往。”\cite{RN2}逊教皆退不用，独唤阶下
末将淳于丹曰：“吾与汝五千军，去取江南第四营：蜀将傅彤所守。今晚就要成功。吾自提
兵接应。”淳于丹引兵去了，又唤徐盛、丁奉曰：“汝等各领兵三千，屯于寨外五里，如淳
于丹败回，有兵赶来，当出救之，却不可追去。”二将自引军去了。

却说淳于丹于黄昏时分，领兵前进，到蜀寨时，已三更之后。丹令众军鼓噪而入。蜀营
内傅彤引军杀出，挺枪直取淳于丹；丹敌不住，拨马便回。忽然喊声大震，一彪军拦住去
路：为首大将赵融。丹夺路而走，折兵大半，正走之间，山后一彪蛮兵拦住：为首番将沙摩
柯。丹死战得脱，背后三路军赶来。比及离营五里，吴军徐盛、丁奉二人两下杀来，蜀兵退
去，救了淳于丹回营。丹带箭入见陆逊请罪。逊曰：“非汝之过也。吾欲试敌人之虚实耳。
破蜀之计，吾已定矣。”徐盛、丁奉曰：“蜀兵势大，难以破之，空自损兵折将耳。”逊笑
曰：“吾这条计，但瞒不过诸葛亮耳。\cite{RN4}


	\end{flushleft}%左对齐结束



	\begin{flushleft}
\heiti{\zihao{3}{
小标题3 议温明董卓叱丁原， 馈金珠李肃说吕布
}}				

\indent\setlength{\parindent}{2em}%首行缩进2字符
						\songti
                  %\zihao{-4}
						% \renewcommand{\baselineskip}{1.5em}
						\hspace{1.6em}%首段首行缩进
天幸此人不在，使我成大功也。”遂集大小将士听
令：使朱然于水路进兵，来日午后东南风大作，用船装载茅草，依计而行；韩当引一军攻江
北岸，周泰引一军攻江南岸，每人手执茅草一把，内藏硫黄焰硝，各带火种，各执枪刀，一
齐而上，但到蜀营，顺风举火；蜀兵四十屯，只烧二十屯，每间一屯烧一屯。各军预带干
粮，不许暂退，昼夜追袭，只擒了刘备方止。众将听了军令，各受计而去。却说先主正在御
营寻思破吴之计，忽见帐前中军旗幡，无风自倒。乃问程畿曰：“此为何兆？”畿曰：“夜
今莫非吴兵来劫营？”先主曰：“昨夜杀尽，安敢再来？”畿曰：“倘是陆逊试敌，奈
何？”正言间，人报山上远远望见吴兵尽沿山望东去了。先主曰：“此是疑兵。”令众休
动，命关兴、张苞各引五百骑出巡。黄昏时分，关兴回奏曰：“江北营中火起。”先主急令
关兴往江北，张苞往江南，探看虚实：“倘吴兵到时，可急回报。”二将领命去了。

初更时分，东南风骤起。只见御营左屯火发。方欲救时，御营右屯又火起。风紧火急，
树木皆着，喊声大震。两屯军马齐出，奔离御营中，御营军自相践踏，死者不知其数。后面
吴兵杀到，又不知多少军马。先主急上马，奔冯习营时，习营中火光连天而起。江南、江
北，照耀如同白日。冯习慌上马引数十骑而走，正逢吴将徐盛军到，敌住厮杀。

先主见了，拨马投西便走。徐盛舍了冯习，引兵追来。先主正慌，前面又一军拦住，乃是吴将丁奉，两
下夹攻。先主大惊，四面无路。忽然喊声大震，一彪军杀入重围，乃是张苞，救了先主，引
御林军奔走。正行之间，前面一军又到，乃蜀将傅彤也，合兵一处而行。背后吴兵追至。先
主前到一山，名马鞍山。张苞、傅彤请先主上的山时，山下喊声又起：陆逊大队人马，将马
鞍山围住。张苞、傅彤死据山口。先主遥望遍野火光不绝，死尸重叠，塞江而下。次日，吴
兵又四下放火烧山，军士乱窜，先主惊慌。忽然火光中一将引数骑杀上山来，视之，乃关兴
也。兴伏地请曰：“四下火光逼近，不可久停。陛下速奔白帝城，再收军马可也。”先主
曰：“谁敢断后？”傅彤奏曰：“臣愿以死当之！”当日黄昏，关兴在前，张苞在中，留傅
彤断后，保着先主，杀下山来。吴兵见先主奔走，皆要争功，各引大军，遮天盖地，往西追
赶，先主令军士尽脱袍铠，塞道而焚，以断后军。正奔走间，喊声大震，吴将朱然引一军从
江岸边杀来，截住去路。先主叫曰：“朕死于此矣！”关兴、张苞纵马冲突，被乱箭射回，
各带重伤，不能杀出。背后喊声又起，陆逊引大军从山谷中杀来。\cite{RN3,RN2} %引用文献RN3


	\end{flushleft}%左对齐结束


	\begin{flushleft}
\heiti{\zihao{3}{
小标题4	废汉帝陈留践位， 谋董贼孟德献刀
}}				

\indent\setlength{\parindent}{2em}%首行缩进2字符
						\songti
                  %\zihao{-4}
						% \renewcommand{\baselineskip}{1.5em}
						\hspace{1.6em}%首段首行缩进
%参考文献插入示范——————————————————————————
杀散吴兵，救出先主，望白帝城而走。\cite{RN2}
%————————————————————————————————
先主曰：“朕虽得脱，诸将士将奈何？”云
曰：“敌军在后，不可久迟。陛下且入白帝城歇息，臣再引兵去救应诸将。”此时先主仅存
百余人入白帝城。后人有诗赞陆逊曰：“持矛举火破连营，玄德穷奔白帝城。一旦威名惊蜀
魏，吴王宁不敬书生。”

    却说傅彤断后，被吴军八面围住。丁奉大叫曰：“川兵死者无数，降者极多，汝主刘备
已被擒获，今汝力穷势孤，何不早降！”傅彤叱曰：“吾乃汉将，安肯降吴狗乎！”挺枪纵
马，率蜀军奋力死战，不下百余合，往来冲突，不能得脱。彤长叹曰：“吾今休矣！”言
讫，口中吐血，死于吴军之中。后人赞傅彤诗曰：“彝陵吴蜀大交兵，陆逊施谋用火焚。至
死犹然骂吴狗，傅彤不愧汉将军。”


	\end{flushleft}%左对齐结束
%\cleardoublepage %这一页留空，余下内容直接进入新的一页
\clearpage{\pagestyle{empty}\cleardoublepage}%空余一页，余下内容直接进入下一页，而且不要出现页眉页码

%—第二章~研究目的以及意义——————————————————————————————

\addcontentsline{toc}{section}{\textbf{\songti{\zihao{-4}{第二章~研究目的以及意义}}}}
	\begin{center}
	\heiti{\zihao{2}{第二章~研究目的以及意义}}
	\end{center}
\vspace*{10pt}

\begin{flushleft}
\indent\setlength{\parindent}{2em}%首行缩进2字符
						\songti
                  %\zihao{-4}
						% \renewcommand{\baselineskip}{1.5em}
						\hspace{1.6em}%首段首行缩进
先主正慌急之间，此时天色已微明，只见前面喊声震天，朱然军纷纷落涧，滚滚投岩：
一彪军杀人，前来救驾。先主大喜，视之，乃常山赵子龙也。时赵云在川中江州，闻吴、蜀
交兵，遂引军出；忽见东南一带火光冲天，云心惊，远远探视，不想先主被困，云奋勇冲杀
而来。陆逊闻是赵云，急令军退。云正杀之间，忽遇朱然，便与交锋；不一合，一枪刺朱然
于马下，杀散吴兵，救出先主，望白帝城而走。

先主曰：“朕虽得脱，诸将士将奈何？”云
曰：“敌军在后，不可久迟。陛下且入白帝城歇息，臣再引兵去救应诸将。”此时先主仅存
百余人入白帝城。后人有诗赞陆逊曰：“持矛举火破连营，玄德穷奔白帝城。一旦威名惊蜀
魏，吴王宁不敬书生。”

    却说傅彤断后，被吴军八面围住。丁奉大叫曰：“川兵死者无数，降者极多，汝主刘备
已被擒获，今汝力穷势孤，何不早降！”傅彤叱曰：“吾乃汉将，安肯降吴狗乎！”挺枪纵
马，率蜀军奋力死战，不下百余合，往来冲突，不能得脱。彤长叹曰：“吾今休矣！”言
讫，口中吐血，死于吴军之中。后人赞傅彤诗曰：“彝陵吴蜀大交兵，陆逊施谋用火焚。至
死犹然骂吴狗，傅彤不愧汉将军。”


	\end{flushleft}%左对齐结束

%\cleardoublepage %这一页留空，余下内容直接进入新的一页
\clearpage{\pagestyle{empty}\cleardoublepage}%空余一页，余下内容直接进入下一页，而且不要出现页眉页码

%—第三章~材料及方法——————————————————————————————
\addcontentsline{toc}{section}{\textbf{\songti{\zihao{-4}{第三章~材料及方法}}}}
	\begin{center}
	\heiti{\zihao{2}{第三章~材料及方法}}
	\end{center}
\vspace*{10pt}

\addcontentsline{toc}{subsection}{\songti{\zihao{-4}{材料}}}
\begin{flushleft}
\heiti{\zihao{2}{
一.材料
}}

\begin{enumerate}

\heiti \zihao{3} \item  入组标准
	
	\begin{enumerate}[(1)] %加入\indent即可让序号变成(1),(2)……缩进一格
\songti
			\zihao{-4}
\item 山坡上马忠引一军出
\item 一箭射中黄忠肩窝
\item 吴兵见忠中箭
\item 一齐来攻
	\end{enumerate}

\heiti \zihao{3} \item  临床资料

\songti \zihao{-4}
\indent\setlength{\parindent}{2em}%首行缩进2字符
\hspace{1.6em}
蜀祭酒程畿，匹马奔至江边，招呼水军赴敌，吴兵随后追来，水军四散奔逃。畿部将叫
曰：“吴兵至矣！程祭酒快走罢！”畿怒曰：“吾自从主上出军，未尝赴敌而逃！”言未
毕，吴兵骤至，四下无路，畿拔剑自刎。

后人有诗赞曰：“慷慨蜀中程祭酒，身留一剑答君王。临危不改平生志，博得声名万古香。”时吴班、张南久围彝陵城，忽冯习到，言蜀兵败，遂引军来救先主，孙桓方才得脱。张、冯二将正行之间，前面吴兵杀来，背后孙桓从彝陵城杀出，两下夹攻。张南、冯习奋力冲突，不能得脱，死于乱军之中。后人有诗赞曰：
“冯习忠无二，张南义少双。沙场甘战死，史册共流芳。”


%插入表格示范———————————————
	\begin{figure}[htbp]%htbp作为浮动体，可以出现在周围文本所在处和一页的顶部
	\centering%内容居中\raggedleft左对齐
	\includegraphics[scale =1.2]{Table_1.pdf}%scaale就是缩放倍数
	\caption*{\zihao{5} \heiti 表1入组病人资料\\  \songti 随便放的一个图，先看看效果。身留一剑答君王。临危不改平生志，博得声名万古香。”时吴班、张南久围彝陵城，忽冯习到，言蜀兵败，遂引军来救先主，孙桓方才得脱。张、冯二将正行之间，前面吴兵杀来，背后孙桓从彝陵城杀出，两下夹攻。张南、冯}%图例figure legend
	\label{Table_1}%用于文章内部自引
	\end{figure}
%——————————————————————



\heiti \zihao{3} \item  细胞株和感受态细胞

\songti \zihao{-4}
\indent\setlength{\parindent}{2em}%首行缩进2字符
吴班杀出重围，又遇吴兵追赶；幸得赵云接着，救回白帝城去了。时有蛮王沙摩柯，匹
马奔走，正逢周泰，战二十余合，被泰所杀。

蜀将杜路，刘宁尽皆降吴。蜀营一应粮草器
仗，尺寸不存。

	\begin{enumerate}[(1)] %加入\indent即可让序号变成(1),(2)……缩进一格
\songti
			\zihao{-4}
\item  	发矫诏诸镇应曹公 
\item  破关兵三英战吕布
\item    	焚金阙董卓行凶 
\item  匿玉玺孙坚背约
	\end{enumerate}



\heiti \zihao{3} \item  试剂与耗材

\songti \zihao{-4}
\indent\setlength{\parindent}{2em}%首行缩进2字符
蜀将川兵，降者无数。时孙夫人在吴，闻猇亭兵败，讹传先主死于军中，遂
驱车至江边，望西遥哭，投江而死。后人立庙江滨，号曰枭姬祠。尚论者作诗叹之曰：“先
主兵归白帝城，夫人闻难独捐生。至今江畔遗碑在，犹著千秋烈女名。”却说陆逊大获全
功，引得胜之兵，往西追袭。

前离夔关不远，逊在马上看见前面临山傍江，一阵杀气，冲天
而起；遂勒马回顾众将曰：“前面必有埋伏，三军不可轻进。”即倒退十余里，于地势空阔
处，排成阵势，以御敌军；即差哨马前去探视。回报并无军屯在此，逊不信，下马登高望
之，杀气复起。逊再令人仔细探视，哨马回报，前面并无一人一骑。逊见日将西沉，杀气越
加，心中犹豫，令心腹人再往探看。

\end{enumerate}


\end{flushleft}%左对齐结束

%__方法_______________________________________________________________________________

\addcontentsline{toc}{subsection}{\songti{\zihao{-4}{方法}}}
\begin{flushleft}
\heiti{\zihao{2}{
二.方法
}}

\begin{enumerate}

\heiti \zihao{3} \item  夷陵之战
	
	\begin{enumerate}[(1)] %加入\indent即可让序号变成(1),(2)……缩进一格
\songti
			\zihao{-4}
\item  	袁绍磐河战公孙 
\item   孙坚跨江击刘表
\item  	王司徒巧使连环计 
\item   董太师大闹凤仪亭
	\end{enumerate}

\heiti \zihao{3} \item  	除凶暴吕布助司徒， 犯长安李傕听贾诩

\songti \zihao{-4}
\indent\setlength{\parindent}{2em}%首行缩进2字符
\hspace{1.6em}
回报江边止有乱石八九十堆，并无人马。逊大疑，令寻
土人问之。须臾，有数人到。逊问曰：“何人将乱石作堆？如何乱石堆中有杀气冲起？”土
人曰：“此处地名鱼腹浦。诸葛亮入川之时，驱兵到此，取石排成阵势于沙滩之上。自此常
常有气如云，从内而起。”陆逊听罢，上马引数十骑来看石阵，立马于山坡之上，但见四面
八方，皆有门有户。逊笑曰：“此乃惑人之术耳，有何益焉！”遂引数骑下山坡来，直入石
阵观看。部将曰：“日暮矣，请都督早回。”

逊方欲出阵，忽然狂风大作，一霎时，飞沙走
石，遮天盖地。但见怪石嵯峨，槎枒似剑；横沙立土，重叠如山；江声浪涌，有如剑鼓之
声。逊大惊曰：“吾中诸葛之计也！”急欲回时，无路可出。正惊疑间，忽见一老人立于马
前，笑曰：“将军欲出此阵乎？”逊曰：“愿长者引出。”老人策杖徐徐而行，径出石阵，
并无所碍，送至山坡之上。逊问曰：“长者何人？”



\heiti \zihao{3} \item  火烧连营

\songti \zihao{-4}
\indent\setlength{\parindent}{2em}%首行缩进2字符
老人答曰：“老夫乃诸葛孔明之岳父黄
承彦也。昔小婿入川之时，于此布下石阵，名八阵图。反复八门，按遁甲休、生、伤、杜、
景、死、惊、开。每日每时，变化无端，可比十万精兵。临去之时，曾分付老夫道：后有东
吴大将迷于阵中，莫要引他出来。老夫适于山岩之上，见将军从死门而入，料想不识此阵，
必为所迷。老夫平生好善，不忍将军陷没于此，故特自生门引出也。”逊曰：“公曾学此阵
法否？”黄承彦曰：“变化无穷，不能学也。”

	\begin{enumerate}[(1)] %加入\indent即可让序号变成(1),(2)……缩进一格
\songti
			\zihao{-4}
\item 山坡上马忠引一军出
\item 一箭射中黄忠肩窝
\item 吴兵见忠中箭
\item 一齐来攻
	\end{enumerate}



\heiti \zihao{3} \item  托孤白帝城

\songti \zihao{-4}
\indent\setlength{\parindent}{2em}%首行缩进2字符
陆逊回寨，叹曰：“孔明真卧龙
也！吾不能及！”于是下令班师。左右曰：“刘备兵败势穷，困守一城，正好乘势击之；今
见石阵而退，何也？”逊曰：“吾非惧石阵而退；吾料魏主曹丕，其奸诈与父无异，今知吾
追赶蜀兵，必乘虚来袭。吾若深入西川，急难退矣。”

遂令一将断后，逊率大军而回。退兵
未及二日，三处人来飞报：“魏兵曹仁出濡须，曹休出洞口，曹真出南郡：三路兵马数十
万，星夜至境，未知何意。”逊笑曰：“不出吾之所料。吾已令兵拒之矣。”正是：雄心方
欲吞西蜀，胜算还须御北朝。未知如何退兵，且看下文分解。

\end{enumerate}

\end{flushleft}%左对齐结束
%\cleardoublepage %这一页留空，余下内容直接进入新的一页
\clearpage{\pagestyle{empty}\cleardoublepage}%空余一页，余下内容直接进入下一页，而且不要出现页眉页码



%—第四章~结果——————————————————————————————
\addcontentsline{toc}{section}{\textbf{\songti{\zihao{-4}{第四章~结果}}}}
	\begin{center}
	\heiti{\zihao{2}{第四章~结果}}
	\end{center}
\vspace*{10pt}
%-结果1——————————————————————-——————————
\addcontentsline{toc}{subsection}{\songti{\zihao{-4}{1.陆逊火烧连营}}}
\begin{flushleft}
\heiti{\zihao{3}{
1.陆逊火烧连营
}}

\songti \zihao{-4}
\indent\setlength{\parindent}{2em}%首行缩进2字符
\hspace{1.6em}
却说章武二年夏六月，东吴陆逊大破蜀兵于猇亭彝陵之地；先主奔回白帝城，赵云引兵
据守。忽马良至，见大军已败，懊悔不及，将孔明之言，奏知先主。先主叹曰：“朕早听丞
相之言，不致今日之败！今有何面目复回成都见群臣乎！”

遂传旨就白帝城住扎，将馆驿改 (\itshape P \upshape = 0.05)
为永安宫。人报冯习、张南、傅彤，程畿、沙摩柯等皆殁于王事，先主伤感不已。又近臣奏
称：“黄权引江北之兵，降魏去了。陛下可将彼家属送有司问罪。”先主曰：“黄权被吴兵
隔断在江北岸，欲归无路，不得已而降魏：是朕负权，非权负朕也，何必罪其家属？”仍给
禄米以养之。却说黄权降魏，诸将引见曹丕，丕曰：“卿今降朕，欲追慕于陈、韩耶？”

%插入Figure（图）的示范——————————————————
%\itshape 意大利斜体 \upshape 正立字体 \quad 空格 
\begin{figure}[htbp]%htbp作为浮动体，可以出现在周围文本所在处和一页的顶部
	\centering%内容居中\raggedleft左对齐
	\includegraphics[scale = 0.7]{Figure_1.pdf}%scaale就是缩放倍数
	\caption*{\zihao{5} \heiti 图1.勤王室马腾举义 报父仇曹操兴师\\  \songti 及刘焉发榜招军时，玄德年已二十八岁矣。当日见了榜文，慨然长叹。随后一人厉声言
曰：“大丈夫不与国家出力，何故长叹？”玄德回视其人，身长八尺，豹头环眼，燕颔虎
须，声若巨雷，势如奔马。玄德见他形貌异常，问其姓名。其人曰：“某姓张名飞，字翼
德。世居涿郡，颇有庄田，卖酒屠猪，专好结交天下豪杰。恰才见公看榜而叹，故此相
问。”玄德曰：“我本汉室宗亲，姓刘，名备。今闻黄巾倡乱，有志欲破贼安民，恨力不
能，故长叹耳。”飞曰：“吾颇有资财，当招募乡勇，与公同举大事，如何。”玄德甚喜，
遂与同入村店中饮酒。}%图例figure legend
	\label{Figure_1}%用于文章内部自引
	\end{figure}
%————————————————————————————

\end{flushleft}

%—结果2—————————————————————————————————————————
\addcontentsline{toc}{subsection}{\songti{\zihao{-4}{2.诸葛亮识破陆逊}}}
\begin{flushleft}
\heiti{\zihao{3}{
2.诸葛亮识破陆逊
}}

\songti \zihao{-4}
\indent\setlength{\parindent}{2em}%首行缩进2字符
\hspace{1.6em}
曹丕问贾诩曰：“朕欲一统天下，先取蜀乎？先取吴乎？”诩曰：“刘备雄才，更兼诸
葛亮善能治国；东吴孙权，能识虚实，陆逊现屯兵于险要，隔江泛湖，皆难卒谋。以臣观
之，诸将之中，皆无孙权、刘备敌手。虽以陛下天威临之，亦未见万全之势也。只可持守，
以待二国之变。”丕曰：“朕已遣三路大兵伐吴，安有不胜之理？”

尚书刘晔曰：“近东吴
陆逊，新破蜀兵七十万，上下齐心，更有江湖之阻，不可卒制，陆逊多谋，必有准备。”丕
曰：“卿前劝朕伐吴，今又谏阻，何也？”晔曰：“时有不同也。昔东吴累败于蜀，其势顿
挫，故可击耳；今既获全胜，锐气百倍，未可攻也。”丕曰：“朕意已决，卿勿复言。”遂
引御林军亲往接应三路兵马。

早有哨马报说东吴已有准备：令吕范引兵拒住曹休，诸葛瑾引
兵在南郡拒住曹真，朱桓引兵当住濡须以拒曹仁。刘晔曰：“既有准备，去恐无益。”丕不
从，引兵而去。”

%\itshape 意大利斜体 \upshape 正立字体 \quad 空格 
\begin{figure}[htbp]%htbp作为浮动体，可以出现在周围文本所在处和一页的顶部
	\centering%内容居中\raggedleft左对齐
	\includegraphics[scale = 0.7]{Figure_2.pdf}%scaale就是缩放倍数
	\caption*{\zihao{5} \heiti 图2.这是图2的小标题\\  \songti 权
泣而奏曰：“臣受蜀帝之恩，殊遇甚厚，令臣督诸军于江北，被陆逊绝断。臣归蜀无路，降
吴不可，故来投陛下。败军之将，免死为幸，安敢追慕于古人耶！”丕大喜，遂拜黄权为镇
南将军。权坚辞不受。忽近臣奏曰：“有细作人自蜀中来，说蜀主将黄权家属尽皆诛戮。”
权曰：“臣与蜀主，推诚相信，知臣本心，必不肯杀臣之家小也。”丕然之。后人有诗责黄
权曰：“降吴不可却降曹，忠义安能事两朝？堪叹黄权惜一死，紫阳书法不轻饶。”}%图例figure legend
	\label{Figure_2}%用于文章内部自引
	\end{figure}
却说吴将朱桓，年方二十七岁，极有胆略，孙权甚爱之；时督军于濡须，闻曹仁引大军
去取羡溪，桓遂尽拨军守把羡溪去了，止留五千骑守城。忽报曹仁令大将常雕同诸葛虔、王
双、引五万精兵飞奔濡须城来。众军皆有惧色。

桓按剑而言曰：“胜负在将，不在兵之多
寡。兵法云：客兵倍而主兵半者，主兵尚能胜于客兵。今曹仁千里跋涉，人马疲困。吾与汝
等共据高城，南临大江，北背山险，以逸待劳，以主制客：此乃百战百胜之势。虽曹丕自
来，尚不足忧，况仁等耶！”于是传令，教众军偃旗息鼓，只作无人守把之状。
\end{flushleft}


%—结果3—————————————————————————————————————————
\addcontentsline{toc}{subsection}{\songti{\zihao{-4}{3.刘备退守白帝城}}}
\begin{flushleft}
\heiti{\zihao{3}{
3.刘备退守白帝城
}}

\songti \zihao{-4}
\indent\setlength{\parindent}{2em}%首行缩进2字符
\hspace{1.6em}
言讫，化阵清风而去。角得此书，晓夜攻习，能呼风唤雨，号
为“太平道人”。中平元年正月内，疫气流行，张角散施符水，为人治病，自称“大贤良
师”。角有徒弟五百余人，云游四方，皆能书符念咒。次后徒众日多，角乃立三十六方，大
方万余人，小方六七千，各立渠帅，称为将军；讹言：“苍天已死，黄天当立；岁在甲子，
天下大吉。”令人各以白土书“甲子”二字于家中大门上。

青、幽、徐、冀、荆、扬、兖、
豫八州之人，家家侍奉大贤良师张角名字。角遣其党马元义，暗赍金帛，结交中涓封谞，以
为内应。角与二弟商议曰：“至难得者，民心也。今民心已顺，若不乘势取天下，诚为可
惜。”遂一面私造黄旗，约期举事；一面使弟子唐周，驰书报封谞。唐周乃径赴省中告变。
帝召大将军何进调兵擒马元义，斩之；次收封谞等一干人下狱。

张角闻知事露，星夜举兵，
自称“天公将军”，张宝称“地公将军”，张梁称“人公将军”。申言于众曰：“今汉运将
终，大圣人出。汝等皆宜顺天从正，以乐太平。”四方百姓，裹黄巾从张角反者四五十万。
贼势浩大，官军望风而靡。何进奏帝火速降诏，令各处备御，讨贼立功。一面遣中郎将卢
植、皇甫嵩、朱儁，各引精兵、分三路讨之。”

且说张角一军，前犯幽州界分。幽州太守刘焉，乃江夏竟陵人氏，汉鲁恭王之后也。当
时闻得贼兵将至，召校尉邹靖计议。靖曰：“贼兵众，我兵寡，明公宜作速招军应敌。”刘
焉然其说，随即出榜招募义兵。


%\itshape 意大利斜体 \upshape 正立字体 \quad 空格 
\begin{figure}[htbp]%htbp作为浮动体，可以出现在周围文本所在处和一页的顶部
	\centering%内容居中\raggedleft左对齐
	\includegraphics[scale = 0.7]{Figure_3.pdf}%scaale就是缩放倍数
	\caption*{\zihao{5} \heiti 图3这是图3的小标题\\  \songti 榜文行到涿县，引出涿县中一个英雄。那人不甚好读书；性宽和，寡言语，喜怒不形于
色；素有大志，专好结交天下豪杰；生得身长七尺五寸，两耳垂肩，双手过膝，目能自顾其
耳，面如冠玉，唇若涂脂；}%图例figure legend
	\label{Figure_3}%用于文章内部自引
	\end{figure}
中山靖王刘胜之后，汉景帝阁下玄孙，姓刘名备，字玄德。昔刘
胜之子刘贞，汉武时封涿鹿亭侯，后坐酎金失侯，因此遗这一枝在涿县。玄德祖刘雄，父刘
弘。弘曾举孝廉，亦尝作吏，早丧。

%插入表格示范———————————————
	\begin{figure}[htbp]%htbp作为浮动体，可以出现在周围文本所在处和一页的顶部
	\centering%内容居中\raggedleft左对齐
	\includegraphics[scale = 1]{Table_2.pdf}%scaale就是缩放倍数
	\caption*{\zihao{5} \heiti 表2.卫健委网上来的资料\\  \songti 正饮间，见一大汉，推着一辆车子，到店门首歇了，入店坐下，便唤酒保：“快斟酒来
吃，我待赶入城去投军。”玄德看其人：身长九尺，髯长二尺；面如重枣，唇若涂脂；丹凤
眼，卧蚕眉，相貌堂堂，威风凛凛。玄德就邀他同坐，叩其姓名。其人曰：“吾姓关名羽，
字长生，后改云长，河东解良人也。因本处势豪倚势凌人，被吾杀了，逃难江湖，五六年
矣。今闻此处招军破贼，特来应募。”玄德遂以己志告之，云长大喜。同到张飞庄上，共议
大事。飞曰：“吾庄后有一桃园，花开正盛；明日当于园中祭告天地，我三人结为兄弟，协
力同心，然后可图大事。”玄德、云长齐声应曰：“如此甚好。”
}%图例figure legend
	\label{Table_2}%用于文章内部自引
	\end{figure}
%——————————————————————

玄德幼孤，事母至孝；家贫，贩屦织席为业。家住本县
楼桑村。其家之东南，有一大桑树，高五丈余，遥望之，童童如车盖。相者云：“此家必出
贵人。”玄德幼时，与乡中小儿戏于树下，曰：“我为天子，当乘此车盖。”叔父刘元起奇
其言，曰：“此儿非常人也！”因见玄德家贫，常资给之。年十五岁，母使游学，尝师事郑
玄、卢植，与公孙瓒等为友。
\end{flushleft}




%—结果4—————————————————————————————————————————
\addcontentsline{toc}{subsection}{\songti{\zihao{-4}{4.诸葛亮临危受命}}}
\begin{flushleft}
\heiti{\zihao{3}{
4.诸葛亮临危受命
}}

\songti \zihao{-4}
\indent\setlength{\parindent}{2em}%首行缩进2字符
\hspace{1.6em}
次日，于桃园中，备下乌牛白马祭礼等项，三人焚香再拜而说誓曰：“念刘备、关羽、
张飞，虽然异姓，既结为兄弟，则同心协力，救困扶危；上报国家，下安黎庶。不求同年同
月同日生，只愿同年同月同日死。皇天后土，实鉴此心，背义忘恩，天人共戮！”誓毕，拜
玄德为兄，关羽次之，张飞为弟。祭罢天地，复宰牛设酒，聚乡中勇士，得三百余人，就桃
园中痛饮一醉。来日收拾军器，但恨无马匹可乘。正思虑间，人报有两个客人，引一伙伴
当，赶一群马，投庄上来。玄德曰：“此天佑我也！”三人出庄迎接。原来二客乃中山大
商：一名张世平，一名苏双，每年往北贩马，近因寇发而回。玄德请二人到庄，置酒管待，
诉说欲讨贼安民之意。二客大喜，愿将良马五十匹相送；又赠金银五百两，镔铁一千斤，以
资器用。


%\itshape 意大利斜体 \upshape 正立字体 \quad 空格 
\begin{figure}[htbp]%htbp作为浮动体，可以出现在周围文本所在处和一页的顶部
	\centering%内容居中\raggedleft左对齐
	\includegraphics[scale = 0.6]{Figure_4.pdf}%scaale就是缩放倍数
	\caption*{\zihao{5} \heiti 图4.这是图4的小标题\\  \songti 玄德谢别二客，便命良匠打造双股剑。云长造青龙偃月刀，又名“冷艳锯”，重八十二
斤。张飞造丈八点钢矛。各置全身铠甲。共聚乡勇五百余人，来见邹靖。邹靖引见太守刘
焉。三人参见毕，各通姓名。玄德说起宗派，刘焉大喜，遂认玄德为侄。不数日，人报黄巾
贼将程远志统兵五万来犯涿郡。刘焉令邹靖引玄德等三人，统兵五百，前去破敌。玄德等欣
然领军前进，直至大兴山下，与贼相见。贼众皆披发，以黄巾抹额。当下两军相对，玄德出
马，左有云长，右有翼德，扬鞭大骂：“反国逆贼，何不早降！”程远志大怒，遣副将邓茂
出战。张飞挺丈八蛇矛直出，手起处，刺中邓茂心窝，翻身落马。程远志见折了邓茂，拍马
舞刀，直取张飞。云长舞动大刀，纵马飞迎。程远志见了，早吃一惊，措手不及，被云长刀
起处，挥为两段。后人有诗赞二人曰：英雄露颖在今朝，一试矛兮一试刀。初出便将威力
展，三分好把姓名标。}%图例figure legend
	\label{Figure_4}%用于文章内部自引
	\end{figure}




众贼见程远志被斩，皆倒戈而走。玄德挥军追赶，投降者不计其数，大胜而回。刘焉亲
自迎接，赏劳军士。次日，接得青州太守龚景牒文，言黄巾贼围城将陷，乞赐救援。刘焉与
玄德商议。玄德曰：“备愿往救之。”刘焉令邹靖将兵五千，同玄德、关、张，投青州来。
贼众见救军至，分兵混战。玄德兵寡不胜，退三十里下寨。

    玄德谓关、张曰：“贼众我寡；必出奇兵，方可取胜。”乃分关公引一千军伏山左，张
飞引一千军伏山右，鸣金为号，齐出接应。次日，玄德与邹靖引军鼓噪而进。贼众迎战，玄
德引军便退。贼众乘势追赶，方过山岭，玄德军中一齐鸣金，左右两军齐出，玄德摩军回身
复杀。三路夹攻，贼众大溃。直赶至青州城下，太守龚景亦率民兵出城助战。贼势大败，剿
戮极多，遂解青州之围。后人有诗赞玄德曰：运筹决算有神功，二虎还须逊一龙。初出便能
垂伟绩，自应分鼎在孤穷。


玄德幼孤，事母至孝；家贫，贩屦织席为业。家住本县
楼桑村。其家之东南，有一大桑树，高五丈余，遥望之，童童如车盖。相者云：“此家必出
贵人。”玄德幼时，与乡中小儿戏于树下，曰：“我为天子，当乘此车盖。”叔父刘元起奇
其言，曰：“此儿非常人也！”因见玄德家贫，常资给之。年十五岁，母使游学，尝师事郑
玄、卢植，与公孙瓒等为友。

\end{flushleft}

%\cleardoublepage %这一页留空，余下内容直接进入新的一页
\clearpage{\pagestyle{empty}\cleardoublepage}%空余一页，余下内容直接进入下一页，而且不要出现页眉页码



%—第五章~讨论——————————————————————————————
\addcontentsline{toc}{section}{\textbf{\songti{\zihao{-4}{第五章~讨论}}}}
	\begin{center}
	\heiti{\zihao{2}{第五章~讨论}}
	\end{center}
\vspace*{10pt}

\songti \zihao{-4}
\hspace{1.6em}
\indent\setlength{\parindent}{2em}%首行缩进2字符

且说魏将先锋常雕，领精兵来取濡须城，遥望城上并无军马。雕催军急进，离城不远，
一声炮响，旌旗齐竖。朱桓横刀飞马而出，直取常雕。战不三合，被桓一刀斩常雕于马下。
吴兵乘势冲杀一阵，魏兵大败，死者无数。朱桓大胜，得了无数旌旗军器战马。曹仁领兵随
后到来，却被吴兵从羡溪杀出。曹仁大败而退，回见魏主，细奏大败之事。丕大惊。正议之
间，忽探马报：“曹真、夏侯尚围了南郡，被陆逊伏兵于内，诸葛瑾伏兵于外，内外夹攻，
因此大败。”言未毕，忽探马又报：”曹休亦被吕范杀败。”丕听知三路兵败，乃喟然叹
曰：“朕不听贾诩、刘晔之言，果有此败！”时值夏天，大疫流行，马步军十死六七，遂引
军回洛阳。吴、魏自此不和。

却说先主在永安宫，染病不起，渐渐沉重，至章武三年夏四日，先主自知病入四肢，又
哭关、张二弟，其病愈深：两目昏花。厌见侍从之人，乃叱退左右，独卧于龙榻之上。忽然
阴风骤起，将灯吹摇，灭而复明，只见灯影之下，二人侍立。

先主怒曰：“朕心绪不宁，教
汝等且退，何故又来！”叱之不退。先主起而视之，上首乃云长，下首乃翼德也。先主大惊
曰：“二弟原来尚在？”云长曰：“臣等非人，乃鬼也。上帝以臣二人平生不失信义，皆敕
命为神。哥哥与兄弟聚会不远矣。”先主扯定大哭。忽然惊觉，二弟不见。即唤从人问之，
时正三更。先主叹曰：“朕不久于人世矣！”遂遣使往成都，请丞相诸葛亮，尚书令李严
等，星夜来永安宫，听受遗命。孔明等与先主次子鲁王刘永、梁王刘理，来永安宫见帝，留
太子刘禅守成都。

且说孔明到永安宫，见先主病危，慌忙拜伏于龙榻之下。先主传旨，请孔
明坐于龙榻之侧。抚其背曰：“朕自得丞相，幸成帝业；何期智识浅陋，不纳丞相之言，自
取其败。悔恨成疾，死在旦夕。嗣子孱弱，不得不以大事相托。”言讫，泪流满面。孔明亦
涕泣曰：“愿陛下善保龙体，以副下天之望！”先主以目遍视，只见马良之弟马谡在傍，先
主令且退。谡退出，先主谓孔明曰：“丞相观马谡之才何如？”孔明曰：“此人亦当世之英
才也。”先主曰：“不然。朕观此人，言过其实，不可大用。丞相宜深察之。”分付毕，传
旨召诸臣入殿，取纸笔写了遗诏，递与孔明而叹曰：“朕不读书，粗知大略。

圣人云：鸟之
将死，其鸣也哀；人之将死，其言也善。朕本待与卿等同灭曹贼，共扶汉室；不幸中道而
别。烦丞相将诏付与太子禅，令勿以为常言。凡事更望丞相教之！”

%\cleardoublepage %这一页留空，余下内容直接进入新的一页
\clearpage{\pagestyle{empty}\cleardoublepage}%空余一页，余下内容直接进入下一页，而且不要出现页眉页码


%—第六章~总结——————————————————————————————
\addcontentsline{toc}{section}{\textbf{\songti{\zihao{-4}{第六章~总结}}}}
	\begin{center}
	\heiti{\zihao{2}{第六章~总结}}
	\end{center}
\vspace*{10pt}

\songti \zihao{-4}
\hspace{1.6em}
\indent\setlength{\parindent}{2em}%首行缩进2字符

孔明曰：“有何圣谕！”先主
泣曰：“君才十倍曹丕，必能安邦定国，终定大事。若嗣子可辅，则辅之；如其不才，君可
自为成都之主。”孔明听毕，汗流遍体，手足失措，泣拜于地曰：“臣安敢不竭股肱之力，
尽忠贞之节，继之以死乎！”言讫，叩头流血。先主又请孔明坐于榻上，唤鲁王刘永、梁王
刘理近前，分付曰：“尔等皆记朕言：朕亡之后，尔兄弟三人，皆以父事丞相，不可怠
慢。”言罢，遂命二王同拜孔明。二王拜毕，孔明曰：“臣虽肝脑涂地，安能报知遇之恩
也！”

先主谓众官曰：“朕已托孤于丞相，令嗣子以父事之。卿等俱不可怠慢，以负朕
望。”又嘱赵云曰：“朕与卿于患难之中，相从到今，不想于此地分别。卿可想朕故交，早
晚看觑吾子，勿负朕言。”云泣拜曰：“臣敢不效犬马之劳！”先主又谓众官曰：“卿等众
官，朕不能一一分嘱，愿皆自爱。”言毕，驾崩，寿六十三岁。时章武三年夏四月二十四日
也。后杜工部有诗叹曰：“蜀主窥吴向三峡，崩年亦在永安宫。翠华想像空山外，玉殿虚无
野寺中。古庙杉松巢水鹤，岁时伏腊走村翁。武侯祠屋长邻近，一体君臣祭祀同。”

先主驾崩，文武官僚，无不哀痛。孔明率众官奉梓宫还成都。太子刘禅出城迎接灵柩，
安于正殿之内。举哀行礼毕，开读遗诏。诏曰：“朕初得疾，但下痢耳；后转生杂病，殆不
自济。朕闻人年五十，不称夭寿。今朕年六十有余，死复何恨？但以卿兄弟为念耳。勉之！
勉之！勿以恶小而为之，勿以善小而不为。惟贤惟德，可以服人；卿父德薄，不足效也。卿
与丞相从事，事之如父，勿怠！勿忘！卿兄弟更求闻达。至嘱！至嘱！”群臣读诏已毕。孔
明曰：“国不可一日无君，请立嗣君，以承汉统。”乃立太子禅即皇帝位，改元建兴。加诸
葛亮为武乡侯，领益州牧。葬先主于惠陵，谥曰昭烈皇帝。尊皇后吴氏为皇太后；谥甘夫人
为昭烈皇后，糜夫人亦追谥为皇后。升赏群臣，大赦天下。

早有魏军探知此事，报入中原。
近臣奏知魏主。曹丕大喜曰：“刘备已亡，朕无忧矣。何不乘其国中无主，起兵伐之？”贾
诩谏曰：“刘备虽亡，必托孤于诸葛亮。亮感备知遇之恩，必倾心竭力，扶持嗣主。陛下不
可仓卒伐之。”正言间，忽一人从班部中奋然而出曰：“不乘此时进兵，更待何时？”众视
之，乃司马懿也。丕大喜，遂问计于懿。懿曰：“若只起中国之兵，急难取胜。须用五路大
兵，四面夹攻，令诸葛亮首尾不能救应，然后可图。”

丕问何五路，懿曰：“可修书一封，
差使往辽东鲜卑国，见国王轲比能，赂以金帛，令起辽西羌兵十万，先从旱路取西平关：此
一路也。再修书遣使赍官诰赏赐，直入南蛮，见蛮王孟获，令起兵十万，攻打益州、永昌、
牂牁、越嶲四郡，以击西川之南：此二路也。

%\cleardoublepage %这一页留空，余下内容直接进入新的一页
\clearpage{\pagestyle{empty}\cleardoublepage}%空余一页，余下内容直接进入下一页，而且不要出现页眉页码



%—参考文献——————————————————————————————
\addcontentsline{toc}{section}{\textbf{\songti{\zihao{-4}{参考文献}}}}
	\begin{center}
	\heiti{\zihao{2}{参考文献}}
	\end{center}
\vspace*{10pt}

\bibliography{NHH.bib}%参考文献从命名为NHH的bib文件中出


%\cleardoublepage %这一页留空，余下内容直接进入新的一页
\clearpage{\pagestyle{empty}\cleardoublepage}%空余一页，余下内容直接进入下一页，而且不要出现页眉页码



%—附录——————————————————————————————
\addcontentsline{toc}{section}{\textbf{\songti{\zihao{-4}{附录}}}}
	\begin{center}
	\heiti{\zihao{2}{附录}}
	\end{center}
\vspace*{10pt}

\addcontentsline{toc}{subsection}{\songti{\zihao{-4}{中英文缩略词表}}}
\begin{center}
\heiti{\zihao{3}{
中英文缩略词表
}}
\end{center}

\clearpage

\addcontentsline{toc}{subsection}{\songti{\zihao{-4}{博士期间发表的论文}}}
\begin{center}
\heiti{\zihao{3}{
博士期间发表的论文
}}
\end{center}
%\cleardoublepage %这一页留空，余下内容直接进入新的一页
\clearpage{\pagestyle{empty}\cleardoublepage}%空余一页，余下内容直接进入下一页，而且不要出现页眉页码



%—致谢——————————————————————————————
\addcontentsline{toc}{section}{\textbf{\songti{\zihao{-4}{致谢}}}}
	\begin{center}
	\heiti{\zihao{2}{致谢}}
	\end{center}
\vspace*{10pt}
\songti \zihao{-4}
\indent\setlength{\parindent}{2em}%首行缩进2字符
\hspace{1.6em}
再遣使入吴修好，许以割地，令孙权起兵十
万，攻两川峡口，径取涪城：此三路也。又可差使至降将孟达处，起上庸兵十万，西攻汉
中：此四路也。然后命大将军曹真为大都督，提兵十万，由京兆径出阳平关取西川；此五路
也。共大兵五十万，五路并进，诸葛亮便有吕望之才，安能当此乎？”

丕大喜，随即密遣能
言官四员为使前去；又命曹真为大都督，领兵十万，径取阳平关。此时张辽等一班旧将，皆
封列侯、俱在冀、徐、青及合淝等处，据守关津隘口，故不复调用。却说蜀汉后主刘禅，自
即位以来，旧臣多有病亡者，不能细说。凡一应朝廷选法，钱粮、词讼等事，皆听诸葛丞相
裁处。时后主未立皇后，孔明与群臣上言曰：“故车骑将军张飞之女甚贤，年十七岁，可纳
为正宫皇后。”后主即纳之。

建兴元年秋八月，忽有边报说：“魏调五路大兵，来取西川；第一路，曹真为大都督，
起兵十万，取阳平关；第二路，乃反将孟达，起上庸兵十万，犯汉中；第三路，乃东吴孙
权，起精兵十万，取峡口入川；第四路，乃蛮王孟获，起蛮兵十万，犯益州四郡；第五路，
乃番王轲比能，起羌兵十万，犯西平关。此五路军马，甚是利害。”已先报知丞相，丞相不
知为何，数日不出视事。后主听罢大惊，即差近侍赍旨，宣召孔明入朝。使命去了半日，回
报：“丞相府下人言，丞相染病不出。”后主转慌；次日，又命黄门侍郎董允、谏议大夫杜
琼，去丞相卧榻前，告此大事。董、杜二人到丞相府前，皆不得入。


\end{document}




















